\subsubsection{Modelo de datos}

El modelo relacional deriva del Formulario 033, de la ficha rápida y de los requerimientos del Objetivo 1. Su entidad central es la atención: cada cita atendida genera una atención, en la que se registran los diagnósticos, el odontograma, los tratamientos y la receta. El modelo completo se consolida en la implementación.

La Tabla~\ref{tab:ob2_entidades} relaciona cada entidad con los requerimientos que atiende. Se añaden dos entidades propias del agente, vinculadas con el paciente y la cita; las guarda la API, porque el servicio del agente no accede a la base de datos, y el estado de la conversación no se almacena.

\begingroup
\fontsize{9}{11}\selectfont
\renewcommand{\arraystretch}{1.2}
\setlength{\LTpre}{0pt}\setlength{\LTpost}{6pt}
\setlength{\tabcolsep}{4pt}
\begin{xltabular}{\linewidth}{|>{\raggedright\arraybackslash}p{3.6cm}|>{\raggedright\arraybackslash}X|>{\raggedright\arraybackslash}p{2.4cm}|}
\caption{Entidades principales del modelo de datos}
\label{tab:ob2_entidades} \\
\hline
\textbf{Entidad} & \textbf{Para qué sirve} & \textbf{Requerimiento} \\ \hline
\endfirsthead
\hline
\textbf{Entidad} & \textbf{Para qué sirve} & \textbf{Requerimiento} \\ \hline
\endhead
Paciente & Identificación y contacto; el teléfono enlaza al paciente con WhatsApp & RF06, RF11 \\ \hline
Historia clínica & Una por paciente; antecedentes del Formulario 033 & RF01 \\ \hline
Cita & Agenda única de recepción y agente; impide cruces por profesional & RF07, RF11, RF12 \\ \hline
Atención & Registro de la consulta en un único momento & RF01, RNF06 \\ \hline
Diagnóstico & Código CIE-10 de la atención & RF03 \\ \hline
Hallazgo del odontograma & Estado de cada pieza y cara: rojo pendiente, azul realizado & RF02 \\ \hline
Tratamiento realizado & Tratamiento del catálogo, con su precio & RF04, RF08 \\ \hline
Receta & Medicamentos con posología predefinida & RF04, RF05 \\ \hline
Pago & Abonos; el saldo se calcula, no se guarda & RF08 \\ \hline
\multicolumn{3}{|l|}{\cellcolor{gray!10}\textbf{Entidades del agente}} \\ \hline
Aceptación del aviso & Teléfono, versión del aviso y fecha de aceptación & RF09 \\ \hline
Recordatorio & Envío del día anterior y respuesta del paciente & RF14 \\ \hline
\end{xltabular}
\endgroup

El modelo sigue cuatro principios:

\begin{itemize}[leftmargin=*]
    \item \textbf{Separación por rol.} Los datos clínicos se guardan aparte de la identificación; solo el odontólogo los edita (RNF01; BN3).
    \item \textbf{Sin borrado físico.} Los registros clínicos y de cobros se anulan, no se eliminan (RNF04; BN2).
    \item \textbf{Valores calculados.} El saldo, la edad y los índices CPO y ceo se calculan, no se almacenan (RF01, RF08).
    \item \textbf{Trazabilidad.} Cada cambio registra el usuario y la fecha (RNF02; BN4).
\end{itemize}

\paragraph{Estados de la cita}

La Figura~\ref{fig:ob2_estados_cita} muestra los estados de la cita y las transiciones permitidas en cada uno.

\begin{figure}[H]
    \centering
    \begin{tikzpicture}[
        font=\footnotesize,
        >={Stealth[length=2mm]},
        estado/.style={rectangle, rounded corners=2pt, draw, minimum width=2.8cm, minimum height=0.9cm, align=center},
        final/.style={estado, fill=gray!15},
        etiqueta/.style={font=\scriptsize, align=center, inner sep=1pt}
    ]
        % Estado inicial
        \node[circle, fill, inner sep=0pt, minimum size=6pt] (ini) at (-4.2,0.9) {};

        % Estados vigentes (transiciones comunes)
        \node[draw, dashed, rounded corners=4pt, minimum width=3.2cm, minimum height=3.2cm, inner sep=0pt] (grupo) at (-1.4,0) {};
        \node[font=\scriptsize] at (-1.4,1.85) {Estados vigentes};
        \node[estado] (pend) at (-1.4,0.9) {PENDIENTE};
        \node[estado] (conf) at (-1.4,-0.9) {CONFIRMADA};

        % Estados finales
        \node[final] (rep) at (6.2,2.7) {REPROGRAMADA};
        \node[final] (noa) at (6.2,0.9) {NO ASISTIÓ};
        \node[final] (ate) at (6.2,-0.9) {ATENDIDA};
        \node[final] (can) at (6.2,-2.7) {CANCELADA};

        % Transiciones
        \draw[->] (ini) -- (pend.west);
        \draw[->] (pend.south) -- node[left, etiqueta] {confirmar} (conf.north);
        \draw[->] (grupo.east) -- node[pos=0.75, sloped, above, etiqueta] {reprogramar\\(sin llegada)} (rep.west);
        \draw[->] (grupo.east) -- node[pos=0.75, sloped, above, etiqueta] {no asistió\\(sin llegada)} (noa.west);
        \draw[->] (grupo.east) -- node[pos=0.75, sloped, below, etiqueta] {atender\\(con llegada)} (ate.west);
        \draw[->] (grupo.east) -- node[pos=0.75, sloped, below, etiqueta] {cancelar} (can.west);
    \end{tikzpicture}
    \caption{Estados de la cita y transiciones según el registro de llegada (en gris, estados finales)}
    \label{fig:ob2_estados_cita}
\end{figure}

La reprogramación crea una nueva cita en estado pendiente, vinculada a la original. Las citas creadas por el agente siguen los mismos estados que las de la recepción.
